\chapter{Введение}
\label{ch:intro}

\section{Введение}

На прошлом занятии мы познакомились с двумя типами модуляции: BPSK и QPSK и программно реализовали их. При реализации логики
PSF нужно было формировать прямоугольный импульс. Я сделал это самым простым методом: повторял значений I и Q L раз. За счёт этого
получал I(t) и Q(t), которые уже длились во времени. Данный подход рабочий и не является ошибкой, но на практике не используется,
поскольку подходит только в случае, когда формирующий фильтр имеет прямоугольную импульсную характеристику. Если импульсная характеристика
имеет форму приподнятого косинуса, то такой метод не сработает. На этом занятии изучим более граммотный подход.

\section{Цель}

Более детально рассмотреть принцип работы формирующего фильтра. Познакомиться с понятием свертки. 
Программно реализовать алгоритм свертки между входящим сигналом и импульсной характеристикой формирующего фильтра. 


\endinput